\section{总结与延伸}

这节课真正展示的是一个 data-to-model control loop：先用 scaling 和 lifecycle cost 决定目标规模，再选择来源与 provenance，接着用 ablation 找 filter，用 governance 处理 opt-out/PII/license，用 formatting/mixture 定义任务，最后用 contamination-aware evaluation 检查模型。StarCoder 的价值不只是 code score，而是这条链可以被审计。

\subsection{端到端因果链}

\begin{importantbox}{从 data decision 到 model behavior}
\begin{enumerate}
  \item \textbf{Source} 决定知识、许可和长尾覆盖；
  \item \textbf{Filter/dedup} 决定质量、diversity 与 memorization；
  \item \textbf{Formatting/mixture} 决定模型看到的任务和上下文；
  \item \textbf{Training scale} 决定容量与 exposure；
  \item \textbf{Evaluation} 决定我们能否区分泛化与 contamination。
\end{enumerate}
任何一层不可见，最终 benchmark 都难以解释。
\end{importantbox}

\subsection{最容易犯的五个错误}

\begin{warningbox}{从 headline 滑向伪结论}
不要把 Chinchilla ratio 写成跨 domain 常数；不要把 stars/comments 当 quality label；不要把 PII detector 写成零风险证明；不要把 HumanEval 分数写成软件工程能力；不要把 open weights 写成完整 openness。课程中的负结果和治理限制与正结果同等重要。
\end{warningbox}

\subsection{复现实验清单}

\begin{lstlisting}[caption={预训练项目的统一复现清单}]
pin_model_tokenizer_optimizer_and_code_commit()
hash_raw_filtered_and_formatted_dataset_versions()
report_source_mixture_repeats_and_effective_tokens()
publish_filter_ablations_seeds_and_negative_results()
document_opt_out_pii_license_and_decontamination()
freeze_evaluation_prompts_harnesses_cutoffs_and_dates()
report_train_and_inference_compute_memory_latency_cost()
\end{lstlisting}

\subsection{自测问题}

\begin{multicols}{2}
\footnotesize
\begin{enumerate}
  \setlength{\itemsep}{0pt}
  \item 为什么 open weights 不等于完整 transparent release？
  \item Kaplan 与 Chinchilla 对 model/data allocation 的差别是什么？
  \item 为什么 compute-optimal 不一定是 lifecycle-optimal？
  \item raw bytes、unique tokens 与 training tokens 有何区别？
  \item Common Crawl 为什么不能直接 tokenize 后训练？
  \item The Stack v1 的 opt-out 需要什么 data lineage？
  \item v2 使用 Software Heritage 带来什么变化？
  \item synthetic data 的 seed、generator、verifier 分别控制什么？
  \item Cosmopedia 怎样增加生成多样性？
  \item 为什么 repository stars filter 训练出最差 ablation？
  \item 多 seeds 对 filter experiment 有什么作用？
  \item MinHash/LSH 为什么适合 near-dedup？
  \item strong dedup 为什么可能提升 code benchmark？
  \item PII detection 与 secret scanning 有何区别？
  \item decontamination 为什么要覆盖 SFT data？
  \item repository/file metadata 怎样改变训练任务？
  \item FIM 如何用 next-token objective 实现？
  \item source sampling weight 怎样改变 effective mixture？
  \item MQA 与 GQA 主要优化什么？
  \item BigCode 的 openness 为什么超过权重发布？
  \item pass@1 与 pass@k 能否直接比较？
  \item membership test 为什么不能证明 memorization？
  \item LiveCodeBench 的 time split 减少哪类泄漏？
  \item 单 GPU 做 code assistant 应优先投资什么？
\end{enumerate}
\end{multicols}

\subsection{拓展阅读}

\begin{multicols}{2}
\footnotesize
\begin{itemize}
  \setlength{\itemsep}{0.15em}
  \item Hoffmann et al., \href{https://arxiv.org/abs/2203.15556}{Training Compute-Optimal Large Language Models}：Chinchilla scaling。
  \item Muennighoff et al., \href{https://arxiv.org/abs/2305.16264}{Scaling Data-Constrained Language Models}：unique data 受限时的重复训练。
  \item Kocetkov et al., \href{https://arxiv.org/abs/2211.15533}{The Stack}；Lozhkov et al., \href{https://arxiv.org/abs/2402.19173}{StarCoder 2 and The Stack v2}。
  \item Penedo et al., \href{https://arxiv.org/abs/2306.01116}{RefinedWeb} 与 \href{https://arxiv.org/abs/2406.17557}{FineWeb}：web filtering 与 ablations。
  \item Lee et al., \href{https://arxiv.org/abs/2107.06499}{Deduplicating Training Data Makes Language Models Better}。
  \item Li et al., \href{https://arxiv.org/abs/2305.06161}{StarCoder: may the source be with you!}：模型、FIM、治理与评测。
  \item Brown et al., \href{https://arxiv.org/abs/2107.03374}{HumanEval}；Jain et al., \href{https://arxiv.org/abs/2403.07974}{LiveCodeBench}：execution 与时间隔离评测。
  \item Zhou et al., \href{https://arxiv.org/abs/2305.11206}{LIMA}；Hu et al., \href{https://arxiv.org/abs/2106.09685}{LoRA}：少量高质量 alignment 与 parameter-efficient fine-tuning。
\end{itemize}
\end{multicols}

\end{document}
